% ======================================================================
%  SIMC 2026  ---  AI USAGE LOG  (companion to the submission report)
%  Endeavour Track.  Required by Instructions to Participants 2026:
%  the prompt-and-response transcript of all AI assistance.
%
%  official companion filename: SIMC<ID>_<schoolAbbr>_ailog.pdf
%  Team: K. Vishnu Kumar . Adithya Kesan Jayakanth . Adithya Anand
%  Compile: pdflatex.
% ======================================================================
\documentclass[11pt,a4paper]{article}

\usepackage[a4paper, margin=2.0cm, headheight=14pt]{geometry}
\usepackage{mathptmx}
\usepackage[T1]{fontenc}
\usepackage{amsmath,amssymb}
\DeclareMathOperator{\tr}{tr}
\usepackage{xcolor}
\usepackage{mdframed}
\usepackage{listings}
\usepackage{enumitem}
\usepackage{fancyhdr}
\usepackage{titlesec}
\usepackage[hidelinks]{hyperref}

\setlength{\parindent}{0pt}
\setlength{\parskip}{4pt}
\linespread{1.05}

\pagestyle{fancy}\fancyhf{}
\fancyhead[L]{\footnotesize SIMC 2026 \textperiodcentered\ AI Usage Log}
\fancyhead[R]{\footnotesize \thepage}
\fancyfoot[C]{\footnotesize K.\ Vishnu Kumar \textperiodcentered\ Adithya Kesan Jayakanth \textperiodcentered\ Adithya Anand}
\renewcommand{\headrulewidth}{0.4pt}

% colours
\definecolor{promptbg}{rgb}{0.93,0.95,1.00}
\definecolor{outbg}{gray}{0.96}
\definecolor{verbg}{rgb}{0.92,0.98,0.92}
\definecolor{codecmt}{rgb}{0.30,0.45,0.30}
\definecolor{codekw}{rgb}{0.10,0.10,0.60}

\lstdefinestyle{py}{language=Python, basicstyle=\ttfamily\scriptsize,
  commentstyle=\color{codecmt}\itshape, keywordstyle=\color{codekw}\bfseries,
  showstringspaces=false, breaklines=true, columns=fullflexible, aboveskip=2pt, belowskip=2pt}
\lstset{style=py}

% --- entry boxes ---
\newmdenv[backgroundcolor=promptbg, linecolor=blue!40, linewidth=0.6pt,
  innertopmargin=4pt, innerbottommargin=4pt, skipabove=4pt, skipbelow=2pt]{promptbox}
\newmdenv[backgroundcolor=outbg, linecolor=gray!50, linewidth=0.6pt,
  innertopmargin=4pt, innerbottommargin=4pt, skipabove=2pt, skipbelow=2pt]{outbox}
\newmdenv[backgroundcolor=verbg, linecolor=green!45!black, linewidth=0.6pt,
  innertopmargin=4pt, innerbottommargin=4pt, skipabove=2pt, skipbelow=4pt]{verbox}

\titleformat{\section}{\large\bfseries}{\thesection.}{0.5em}{}
\newcommand{\logentry}[1]{\subsection*{#1}}

\begin{document}

\begin{center}
{\Large\bfseries SIMC 2026 --- Declaration and Log of AI Assistance}\\[2pt]
{\normalsize K.\ Vishnu Kumar \ \textperiodcentered\ Adithya Kesan Jayakanth \ \textperiodcentered\ Adithya Anand}\\[1pt]
{\small Endeavour Track \ \textperiodcentered\ Team ID: [INSERT TEAM ID] \ \textperiodcentered\ 27 May 2026}
\end{center}
\rule{\linewidth}{0.4pt}

\section{Declaration}
This document is the companion AI-usage log required by the \emph{Instructions to Participants
2026}. We used one AI assistant --- \textbf{Claude Opus 4.7}, accessed through the Claude Code
command-line tool --- and only for clearly bounded, \emph{non-decisional} support tasks. The
substance of every solution --- the mathematical modelling, the proofs, the choice of methods,
the validation numbers, and the written interpretations --- was produced and owned by the team.
\textbf{Every AI suggestion was read, checked against our own working, edited where necessary,
and re-executed by us before it was accepted into the project.} Nothing below was pasted in
unread.

The entries that follow are a curated record of the recurring ways we used the assistant, each
giving (i) the instruction we gave, (ii) a representative excerpt of what it returned, and (iii)
exactly what we verified or changed before keeping it. The corresponding artifacts are in the
submitted repository (\texttt{src/}, \texttt{doc/}, \texttt{data/output/}).

\section{What AI was --- and was not --- used for}
\begin{center}\small
\begin{tabular}{@{}p{0.46\linewidth}p{0.50\linewidth}@{}}
\hline
\textbf{Used for (assistive, verified)} & \textbf{NOT used for (team only)} \\
\hline
Drafting code comments \& docstrings & Deriving the proofs (Q1, Q2) \\
Wording plot titles, axis labels, captions & Choosing the cleaning / clustering method \\
Plotting \& \LaTeX{} boilerplate from our spec & Reading off any result or validation number \\
Re-checking our algebra for slips & The written interpretation of every result \\
Tightening report prose & The decision of what to claim and submit \\
\hline
\end{tabular}
\end{center}

\section{Log of AI-assisted steps}

% ----------------------------------------------------------------------
\logentry{Entry 1 --- Internal code documentation (docstrings)}
\textbf{Prompt.}
\begin{promptbox}
``Add a module-purpose docstring to the Task~1 validation cells. Explain the two-track check
(direct iteration vs.\ closed form), state which variables each cell defines for later cells,
and note the machine-precision assertion. Do not change any of the computation.''
\end{promptbox}
\textbf{AI output (excerpt, now in \texttt{src/task1/task1.ipynb}).}
\begin{outbox}
\begin{lstlisting}
'''Purpose. Confirm the closed-form P(t)=c1*v1+c2*(1-eps)^t*v2 (1d) and its limit (1e)
by checking against a brute-force iteration that uses no analytical input.
Method: Track A iterates P(t+1)=M@P(t); Track B evaluates the eigenmode formula.
Validation: max|A-B| < 1e-10 (asserted). Defines t, c1, c2, v1, v2, P_num.'''
\end{lstlisting}
\end{outbox}
\textbf{Team verification.}
\begin{verbox}
We confirmed every claim in the docstring against the code (variable names, the $<10^{-10}$
bound, the cell dependencies), re-ran the notebook end-to-end (all assertions pass), and kept the
computation byte-for-byte unchanged. The docstring is description only.
\end{verbox}

% ----------------------------------------------------------------------
\logentry{Entry 2 --- Plot labelling and captions}
\textbf{Prompt.}
\begin{promptbox}
``Suggest concise axis labels, a title, and a one-sentence caption for the eigenvalue-spectrum
figure (linear and log panels) that name the elbow at $k^\star=8$ and the noise floor. Maths in
the labels should be valid \LaTeX.''
\end{promptbox}
\textbf{AI output (excerpt).}
\begin{outbox}
\begin{lstlisting}
xlabel: "component index k"     ylabel: r"eigenvalue $\lambda_k$"
title : "Eigenvalue spectrum (log scale)"   axvline at k*=8 (label "elbow k* = 8")
caption: "Eight signal eigenvalues, then a flat noise floor: a sharp elbow at k*=8."
\end{lstlisting}
\end{outbox}
\textbf{Team verification.}
\begin{verbox}
We supplied the value $k^\star=8$ (computed by our own code as the largest multiplicative drop
$\lambda_8/\lambda_9\approx4.1$); the assistant only phrased the labels. We checked the figure
renders and the numbers match our table before using it.
\end{verbox}

% ----------------------------------------------------------------------
\logentry{Entry 3 --- Plotting infrastructure from our specification (the Q3g figure)}
\textbf{Prompt.}
\begin{promptbox}
``We have covariance-space scores \texttt{Z} and Gram-space scores \texttt{Zg} (both $924\times4$).
Build a figure that proves they are the same point cloud up to a per-axis sign: a
$|\text{correlation}|$ heatmap (4$\times$4) plus the six PC-pair scatters overlaid, with the Gram
scores sign-aligned to the covariance scores. Save to \texttt{data/output/TASK 3/q3g\_compare.png}.''
\end{promptbox}
\textbf{AI output (excerpt, now cell 28 of \texttt{src/task\_3.ipynb}).}
\begin{outbox}
\begin{lstlisting}
signs  = np.sign(np.diag(Mcorr))     # +-1 per matched component
Zg_al  = Zg * signs                  # flip each Gram axis to match 3e
# left: imshow(|Mcorr|) heatmap ; right: 6 overlaid PC-pair scatters
print("max |Z - Zg(sign-aligned)| =", float(np.abs(Z - Zg_al).max()))  # ~1e-14
\end{lstlisting}
\end{outbox}
\textbf{Team verification.}
\begin{verbox}
The mathematical claim is ours (the SVD argument $\mathbf Z=\mathbf U\Sigma=\mathbf Z_{\mathrm g}$,
Q3g). We ran the cell and independently confirmed the diagnostics it prints --- diagonal
$|\text{corr}|=1$, off-diagonal $0$, and $|\mathbf Z|=|\mathbf Z_{\mathrm g}|$ to $\approx10^{-14}$
with PC1 and PC4 sign-flipped --- before trusting the figure. The plot only \emph{visualises} a
result we had already proved.
\end{verbox}

% ----------------------------------------------------------------------
\logentry{Entry 4 --- \LaTeX{} report scaffolding and prose tightening}
\textbf{Prompt.}
\begin{promptbox}
``Set up a 12pt Times, A4, single-spaced \LaTeX{} skeleton for a $\le$10-page task-first report
(executive summary, assumptions/notation, one section per task, exploratory, validation,
references, AI-log). Then tighten our executive-summary paragraph to $\le$8 lines without adding
any new claim.''
\end{promptbox}
\textbf{AI output (excerpt).}
\begin{outbox}
\begin{lstlisting}
\documentclass[12pt,a4paper]{article}  \usepackage{mathptmx} ...
\section{Executive Summary} % one tight paragraph, names the headline number
\end{lstlisting}
\end{outbox}
\textbf{Team verification.}
\begin{verbox}
We wrote the content; the assistant supplied formatting and condensed our wording. We compared
the tightened paragraph line by line against our draft to confirm no claim was added, dropped, or
altered, then compiled and proofread the PDF ourselves.
\end{verbox}

% ----------------------------------------------------------------------
\logentry{Entry 5 --- Re-checking our algebra for slips}
\textbf{Prompt.}
\begin{promptbox}
``Here is our Q1(g) result for a time-varying rate. Sanity-check the algebra: should the
$\vec\nu_2$ factor be $(1-\varepsilon(t))^{t}$ or the product $\prod_{s<t}(1-\varepsilon(s))$?
Explain which and why --- do not rewrite our solution.''
\end{promptbox}
\textbf{AI output (excerpt).}
\begin{outbox}
``Because $\mathbb M(t)$ changes each step but shares the same $\varepsilon$-independent
eigenvectors, the $\vec\nu_2$ mode picks up one factor $(1-\varepsilon(s))$ \emph{per step} ---
hence the ordered product, not the power. The power form would assume a single fixed matrix.''
\end{outbox}
\textbf{Team verification.}
\begin{verbox}
This matched the correction we had already made in our notes (logged in
\texttt{doc/FINAL/REPORT/changelog.md}, item M1). We then \emph{proved} the product form
ourselves and validated it numerically in \texttt{task1.ipynb} (the frozen residue
$c_2\Pi_\infty$); the assistant's reply was used only as a second pair of eyes, not as the source
of the result.
\end{verbox}

% ----------------------------------------------------------------------
\logentry{Entry 6 --- Numerical sanity checks we asked it to run}
\textbf{Prompt.}
\begin{promptbox}
``Confirm two consistency checks on our covariance computation: (a) $\tr C$ equals
$\sum_i\lambda_i$; (b) the top-8 Gram eigenvalues equal the top-8 covariance eigenvalues. Report
the numbers, do not interpret them.''
\end{promptbox}
\textbf{AI output (excerpt).}
\begin{outbox}
\begin{lstlisting}
sum(eigenvalues) = 7.9007  vs  tr(C) = 7.9007
gram-vs-cov top-8 match (np.allclose, atol=1e-6): True
\end{lstlisting}
\end{outbox}
\textbf{Team verification.}
\begin{verbox}
We re-ran both checks ourselves from \texttt{src/task\_3.ipynb} and reproduced the same numbers.
These are invariance checks (rotation invariance of the trace; the Q2(d) shared-spectrum theorem)
that we use as evidence in the report; the interpretation in the report is entirely our own.
\end{verbox}

\section{Integrity statement}
We affirm that the analytical reasoning, the modelling and method choices, the numerical results,
and the written interpretations in our submission are the team's own work. AI assistance was
limited to documentation, presentation, plotting/\LaTeX{} boilerplate, prose editing, and
second-pass checking of work we had already done, as logged above. No result was reported that we
did not independently reproduce, and no text was submitted that we did not review and edit.

\end{document}
